%% kapa: exam=1 round=22 year=2011 period=2 form=A subject=accounting subject_ko=회계학 seq=37 number=37
%% source: Q-Net 감정평가사 기출문제 (cst00309) artlSeq=5036790 file=extracted/1008/5036790/2011년 제22회 감정평가사 1차시험 기출문제/2교시.hwp
%% conversion: kordoc 4.24.51 -> tools/gen1.py (auto); review=proofread-sonnet; flags=none
\begin{question}{37}
20×1년에 설립된 (주)서울은 제1공정에서 원재료 1,000kg을 가공하여 중간제품A와 제품B를 생산한다. 제품B는 분리점에서 즉시 판매될 수 있으나, 중간제품A는 분리점에서 판매가치가 형성되어 있지 않기 때문에 제2공정에서 추가 가공하여 제품C로 판매한다. 제품별 생산 및 판매량과 kg당 판매가격은 다음과 같다.

{\small\centering
\begin{tabular}{|l|r|r|}
\hline
\multicolumn{1}{|c|}{제 품} & \multicolumn{1}{c|}{생산 및 판매량} & \multicolumn{1}{c|}{kg당 판매가격} \\
\hline
중간제품A & 600kg & - \\
제품B & 400kg & ￦500 \\
제품C & 600kg & 450 \\
\hline
\end{tabular}\par}

제1공정에서 발생한 결합원가는 ￦1,200,000이었고, 중간제품A를 제품C로 가공하는데 추가된 원가는 ￦170,000이었다. 회사가 결합원가를 순실현가치에 비례하여 제품에 배부하는 경우, 제품B와 제품C에 배부되는 총제조원가는? \par
\noindent\hspace*{2em}\makebox[2.8cm][l]{\uline{제품B}}\uline{제품C}\par
\begin{choices}
  \item \makebox[2.8cm][l]{￦400,000}￦800,000
  \item \makebox[2.8cm][l]{￦400,000}￦970,000
  \item \makebox[2.8cm][l]{￦570,000}￦800,000
  \item \makebox[2.8cm][l]{￦800,000}￦570,000
  \item \makebox[2.8cm][l]{￦870,000}￦400,000
\end{choices}
\end{question}
\answer{4}
